\documentclass[11pt,a4paper]{article}

\usepackage{graphicx}
\usepackage{xcolor}
\usepackage[left=2.5cm, right=2.5cm, top=2.5cm, bottom=2.5cm]{geometry}
\usepackage{amsmath,amssymb,amsthm}
\usepackage{fontspec}
\setmainfont{Liberation Serif}
\setsansfont{Liberation Sans}
\setmonofont{Liberation Mono}[Scale=0.86]
\usepackage{booktabs}
\usepackage{tabularx}
\newcolumntype{Y}{>{\raggedright\arraybackslash}X}
\usepackage{colortbl}
\usepackage{adjustbox}
\usepackage{enumitem}
\usepackage[htt]{hyphenat}
\usepackage{tcolorbox}
\definecolor{linkblue}{HTML}{1e3a5f}
\definecolor{accgreen}{HTML}{0e7c66}
\definecolor{warnamber}{HTML}{b45309}
\definecolor{rowlight}{HTML}{f0f5fa}
\usepackage[
    colorlinks=true,
    linkcolor=linkblue,
    citecolor=darkgray,
    urlcolor=linkblue,
    bookmarks=true,
    bookmarksnumbered=true,
    unicode=true
]{hyperref}

\widowpenalty=10000
\clubpenalty=10000
\emergencystretch=1.8em
\tolerance=2500
\hbadness=2500
\newcommand{\dd}{\mathrm{d}}

\newtheorem{theorem}{Theorem}[section]
\newtheorem{lemma}[theorem]{Lemma}
\newtheorem{proposition}[theorem]{Proposition}
\newtheorem{corollary}[theorem]{Corollary}
\theoremstyle{definition}
\newtheorem{definition}[theorem]{Definition}
\theoremstyle{remark}
\newtheorem{remark}[theorem]{Remark}

\newtcolorbox{honestybox}[1][]{colback=warnamber!6, colframe=warnamber!80,
  title={\textbf{Honesty note:} #1}, fonttitle=\small, boxrule=0.7pt, arc=2pt}
\newtcolorbox{findbox}[1][]{colback=accgreen!7, colframe=accgreen!80,
  title={\textbf{Finding:} #1}, fonttitle=\small, boxrule=0.7pt, arc=2pt}

\hypersetup{
    pdftitle={The Choptuik Problem: Statement, Solvability, the Attack on 224 = 4·56, and the e >= 1 Threshold},
    pdfauthor={Super Z (Z.ai)},
    pdfsubject={Klein quartic, the {7,3} tiling, Cuntz defect, critical collapse},
    pdfkeywords={Choptuik, Klein quartic, 7-3 tiling, PSL(2,7), Cuntz defect, partial isometry, DSI-4}
}

\begin{document}

\tableofcontents
\newpage

% ==========================================================================
\section{Summary of results}\label{sec:summary}
% ==========================================================================

This report continues the computational cycle on the Choptuik--Isaev framework
(the repository \texttt{choptuik\_ac\_bc}) and answers three questions posed
for this session: what exactly is the Choptuik problem and can it be solved;
why does the DSI-4 observation single out the denominator $224 = 4\cdot 56$,
a constant tied to the combinatorics of the $\{7,3\}$ tiling of the Klein
quartic; and whether the ``obstruction ladder'' of the Cuntz defect can be
turned into a theorem by proving the threshold $e \ge 1$ for constructions
isometric on their carrier. All computations were carried out by two new
deterministic experiments of the package \texttt{choptuik\_experiments/},
namely \texttt{exp6\_census\_224.py} and \texttt{exp7\_e1\_threshold.py};
generator seeds are fixed and every run is reproduced by a single command.

The main results are as follows. First, for the class of
\emph{carrier-isometric} pairs --- pairs of partial isometries
$X_1, X_2 \in M_n(\mathbb{C})$, which contains all ``structural'' models of
the framework (permutation matrices, truncated Cuntz-tree shifts) --- we
prove the threshold $e = F(X_1,X_2)/n \ge 1$: the value $5/7$ that minimises
the defect over all matrix pairs is unattainable in this class because it
would require contracting the carrier. Together with the exact covering
formula $F = n + 2\,\mathrm{rank}(I - P - Q)$ and the identity $F = 2n$ for
unitary pairs, this turns the ladder $5/7 \to 1 \to 2$ into a theorem with
an exact defect arithmetic (Section~\ref{sec:threshold}). Second, the attack
on DSI-4 shows that $224 = 4\cdot 56 = 8\cdot 28$ is a \emph{census}
constant of the $\{7,3\}$ tiling: it is expressible through the Klein quartic
census in at least seven independent ways, and the fraction $9/224$ is the
\emph{unique} census fraction within $0.1\%$ of the measured $c_{K3}$
(Section~\ref{sec:attack224}). The main acquisition is the genus family
$c(g)$: the two structural readings of $27/672$ coincide \emph{only} at
$g = 3$, so a single measurement of the DSI amplitude on any higher Hurwitz
surface discriminates the mechanisms. Third, the Choptuik problem statement
is fixed at both levels (physical and spectral), and its solvability is
assessed at three levels of rigour: numerically the problem is solved, the
rigorous theory is partial, and a closed form is unknown
(Sections~\ref{sec:statement} and~\ref{sec:solvability}).

\begin{table}[htbp]
\centering
\small
\caption{Session outcomes: results and their status.}
\label{tab:summary}
\begin{tabularx}{\textwidth}{@{}lYY@{}}
\toprule
Result & Statement & Status \\
\midrule
Theorem~A (\S\ref{sec:threshold}) & $e \ge 1$ for partial isometries;
equality $\Leftrightarrow$ $PQ = 0$, $p + q = n$ & proved + verified
(0 violations in 24{,}000 pairs) \\
Theorem~B (\S\ref{sec:threshold}) & $F = n + 2\,\mathrm{rank}(I{-}P{-}Q)$
exactly; tree: $F = n + 2$ & proved (accuracy $10^{-14}$) \\
Theorem~C (\S\ref{sec:threshold}) & unitary pairs: $F = 2n$ exactly &
proved (accuracy $10^{-14}$) \\
Uniqueness lemma (\S\ref{sec:attack224}) & $9/224$ is the unique census
fraction within $0.1\%$ of $c_{K3}$ & computational, exhaustive search \\
Degeneracy theorem (\S\ref{sec:attack224}) & $c_I(g) = c_{II}(g)
\Leftrightarrow g = 3$ & proved (algebra + grid $g \le 100$) \\
Discriminator (\S\ref{sec:attack224}) & branches $(g{-}1)^{-1}$ vs
$(g{-}1)^{-2}$ separable by one measurement at $g \ge 7$ &
falsifiable prediction \\
\bottomrule
\end{tabularx}
\end{table}

% ==========================================================================
\section{The Choptuik problem: statement}\label{sec:statement}
% ==========================================================================

\subsection{Physical statement: critical collapse of a scalar field}

The physical Choptuik problem (Choptuik, 1993) is gravitational collapse of
a massless scalar field in spherical symmetry. One studies the
Einstein--Klein--Gordon system $G_{\mu\nu} = 8\pi\,T_{\mu\nu}[\phi]$ with
$T_{\mu\nu} = \partial_\mu\phi\,\partial_\nu\phi
- \tfrac{1}{2}(\partial\phi)^2 g_{\mu\nu}$ and the wave equation
$\Box\phi = 0$; in the radial gauge used by the monograph (the Tyukovsky
formalism) the metric is
\begin{equation}\label{eq:metric}
\dd s^2 = -\alpha(r,t)^2\,\dd t^2 + a(r,t)^2\,\dd r^2 + r^2\,\dd\Omega^2,
\end{equation}
and the field--metric equations reduce to a coupled system for the three
functions $(\phi, a, \alpha)$: the scalar equation
$\tfrac{1}{a}\partial_r(\tfrac{1}{a}\partial_r\phi)
- \tfrac{1}{\alpha}\partial_t(\tfrac{1}{\alpha}\partial_t\phi) = V'(\phi)$
with two constraints on $\partial_t a$ and $\partial_r\alpha$
(equations (9.3)--(9.5) of the monograph). A one-parameter family of smooth
initial data is labelled by the amplitude $p$; a critical value $p_*$
separates dispersion ($M_{BH} = 0$) from collapse. Near $p_*$ Choptuik
discovered three universal phenomena. First, \emph{mass scaling}:
$M_{BH} \propto (p - p_*)^{\gamma}$ with an exponent $\gamma \approx 0.374$
independent of the data family. Second, \emph{discrete self-similarity}
(DSS): the critical solution reproduces itself under logarithmic scale
translation with echoing period $\Delta \approx 3.44$. Third, universality
of the critical solution itself. The mechanism is explained by perturbation
theory around the DSS solution: there is exactly one unstable mode with
growth exponent $\lambda_0 \approx 2.674$, and $\gamma = 1/\lambda_0$;
the scaling is illustrated in Fig.~\ref{fig:choptuik}.

\begin{figure}[htbp]
\centering
\includegraphics[width=0.98\textwidth]{figs_en/fig_choptuik.png}
\caption{Left: universal mass scaling $M_{BH} \propto (p-p_*)^\gamma$ for two
distinct families of initial data (synthetic illustration of the 1993
statement). Right: the measured exponent $\gamma \approx 0.374$ versus the
framework candidate $b_{\mathrm{Ch}} = 1 - \cos(2\pi/7) \approx 0.3765$;
the gap is $+0.67\%$.}
\label{fig:choptuik}
\end{figure}

\subsection{The ``Choptuik problem'' of the monograph: the spectral gap of the Klein quartic}

In the Isaev monograph, the ``Choptuik problem'' is a related but distinct
statement: the \emph{exact determination of the spectral gap}
$\Delta = 3.443$ of the Laplacian on the Klein quartic
$K\colon x^3y + y^3z + z^3x = 0 \subset \mathbb{CP}^2$ (genus $g = 3$,
$\mathrm{Aut}(K) \cong \mathrm{PSL}(2,7)$, $|\mathrm{PSL}(2,7)| = 168$)
equipped with the canonical hyperbolic metric. The uniformisation is given
by the Fuchsian group $\Gamma(2,3,7) = \langle A,B,C \mid A^2 = B^3 = C^7
= ABC = I \rangle$; the spinor phases of the generators $\delta_A = \pi/2$,
$\delta_B = \pi/3$, $\delta_C = \pi/7$ produce two corrections to the gap,
and the unified formula of the monograph,
\begin{equation}\label{eq:delta_ch}
\Delta_{\mathrm{Ch}} \;=\; \lambda_1 - \frac{R}{4} + \frac{\delta_C^2}{2}
- \frac{\delta_C^5}{22} \;=\; 3.338 + \frac{\delta_C^2}{2} - \frac{\delta_C^5}{22}
\;\approx\; 3.4379,
\end{equation}
reproduces the observed $\Delta = 3.443$ within $0.15\%$. The same machinery
carries the critical-exponent constant of the framework,
\begin{equation}\label{eq:bch}
b_{\mathrm{Ch}}(n) = 1 - \cos\frac{2\pi}{n},
\qquad b_{\mathrm{Ch}}(7) = 0.37651\ldots,
\end{equation}
identified (to $0.67\%$) with the physical $\gamma \approx 0.374$; the DSS
echoing period $\Delta \approx 3.44$ additionally resonates with the gap
$3.443$ (mismatch $\sim 0.1\%$). In the broad sense, the ``Choptuik problem''
is thus the programme of reducing the three numbers of critical collapse
($\gamma$, $\Delta$, and the DSI-modulation amplitude
$c_{K3} = 0.040177576\ldots$) to the spectrum and combinatorics of the
genus-3 arithmetic curve.

\begin{honestybox}[status of the bridge]
At present the connection between the Klein quartic and collapse is an
analogy resting on numerical coincidences at the $0.1\%$--$0.8\%$ level: the
gap $\Delta_{\mathrm{Ch}}$ (formula~\eqref{eq:delta_ch}), the exponent
$b_{\mathrm{Ch}}(7)$, and the echoing period. A mechanism that would derive
the perturbation spectrum of the DSS metric from the Maass spectrum of the
arithmetic curve is absent; Section~\ref{sec:solvability} records this as
the central gap of the programme.
\end{honestybox}

% ==========================================================================
\section{Can the Choptuik problem be solved?}\label{sec:solvability}
% ==========================================================================

The answer depends on the level of ambition; three levels are useful.

\textbf{Level I: numerical solution.} Here the problem \emph{is solved} and
continues to be refined by routine means. Adaptive-grid simulations combined
with perturbation theory around the DSS solution give $\gamma = 0.374(1)$,
echoing period $\Delta \approx 3.44$, and confirm exactly one unstable mode;
further progress is a matter of computational engineering, not principle.
For the framework this level sets a hard bar: any analytic ``value'' of the
exponent must agree with $\gamma$ to $\sim 10^{-3}$, otherwise it is
falsified. The candidate $b_{\mathrm{Ch}}(7) = 0.3765$ passes this bar only
marginally: the $+0.67\%$ gap is five times the modern uncertainty, and as
an \emph{exact} claim it is ruled out unless the $\delta$-machinery of the
monograph (the terms $\delta_C^2/2$, $\delta_C^5/22$ and higher) derives a
correction bringing the value to $0.374$. Note that for the gap the
analogous programme has been carried through: formula~\eqref{eq:delta_ch}
achieves $0.15\%$, with the residue attributed to higher-order holonomy
terms.

\textbf{Level II: rigorous mathematical theory.} Here the picture is partial.
For the spherically symmetric scalar field, Christodoulou's work
(1991--1999) establishes the global evolution picture: dispersion, black-hole
formation, instability of naked singularities, and weak cosmic censorship.
However, the \emph{discrete} self-similarity of the critical solution resists
rigour: existence and uniqueness of the DSS critical solution as theorems,
a proof of universality, and a rigorous count of unstable modes all remain
open problems of mathematical general relativity. Bounded-variation
techniques, which work for CSS-like problems, do not transfer directly to
the oscillating DSS regime; the self-similar change of variables makes the
equation coefficients singular in the echoing period. Solving the Choptuik
problem in the strict sense is therefore a programme at the level of the
best current problems of mathematical GR, not of a single report.

\textbf{Level III: a closed form for the exponent.} No closed formula for
$\gamma$ is known; moreover, there is no substantive reason to expect one:
the exponent is the reciprocal growth rate of the single unstable eigenvalue
of the linearised operator on a DSS background, and eigenvalues of such
operators are typically transcendental. The framework's programme claims the
opposite --- that $\gamma$ is expressible through the arithmetic of genus 3.
For this claim to become a theory, two things are needed. First, a
\emph{mechanism}: an explicit correspondence between the DSS perturbation
problem and objects of the quartic (Maass spectrum, the $\{7,3\}$ tiling,
spin structures). Second, \emph{precision}: an explanation of the $0.67\%$
gap between $b_{\mathrm{Ch}}(7)$ and $\gamma$ that is not a fit. Our
contribution here is twofold: the theorems of Section~\ref{sec:threshold}
close an entire class of candidate ``mechanisms'' (isometric transports),
and the genus family of Section~\ref{sec:attack224} turns the single
coincidence $27/672$ into a falsifiable one-parameter family of predictions
--- that is, into testable physics rather than numerology.

\begin{findbox}[solvability verdict]
Numerically the Choptuik problem is solved ($\gamma = 0.374(1)$,
$\Delta \approx 3.44$, one unstable mode). A rigorous theory of the DSS
critical solution is a major open problem of mathematical GR. A closed form
of $\gamma$ via Klein-quartic combinatorics is an unproven programme of the
framework with two precise gaps --- the mechanism and the $0.67\%$ offset ---
both now formulated with an attack plan.
\end{findbox}

% ==========================================================================
\section{Analytical attack on DSI-4: why $224 = 4\cdot 56$?}\label{sec:attack224}
% ==========================================================================

The DSI-4 observation records the log-periodic modulation amplitude
$c_{K3}$, whose extended writing equals $0.04017757639214903$, and its
integer reading $c_{K3} \approx 27/672 = 9/224$ (deviation $+0.0025\%$).
This section asks: does the combinatorics of the $\{7,3\}$ tiling derive
the denominator $224 = 4\cdot 56$ naturally and uniquely? The answer: yes
--- the constant is \emph{census-born}, and its census status is now
established constructively.

\subsection{The census of the $\{7,3\}$ tiling from $G = \mathrm{GL}(3,2)$}

We constructed the tiling combinatorics explicitly from the group
$G = \mathrm{GL}(3,2) \cong \mathrm{PSL}(2,7)$ of 168 matrices $3\times 3$
over $\mathbb{F}_2$. The element census by order gives $1, 21, 56, 42, 48$
elements of orders $1, 2, 3, 4, 7$ respectively; the Sylow numbers are
$n_2 = 21$, $n_3 = 28$, $n_7 = 8$. The cells of the tiling are cosets by the
stabilisers: vertices $G/H_3$ ($H_3 \cong C_3$; three heptagons meet at each
vertex), edges $G/H_2$, faces $G/H_7$; this yields $V = 56$, $E = 84$,
$F = 24$ and the Euler characteristic $V - E + F = -4 = 2 - 2g$. The
integrity of the construction was verified by two nontrivial graphs: the
1-skeleton of the tiling (the Klein graph) is a connected cubic graph on 56
vertices with 84 edges and girth 7 (such a graph is unique), and the
face-adjacency graph is a 7-valent graph on 24 vertices with legitimate
triangles (three heptagons at each vertex are pairwise adjacent). The three
valence streams coincide: $7F = 3V = 2E = |G| = 168$ --- the valence
identity of the regular $\{7,3\}$ tiling of genus 3.

\begin{table}[htbp]
\centering
\small
\caption{Census of the Klein quartic and the $\{7,3\}$ tiling (all values
derived by the explicit $\mathrm{GL}(3,2)$ construction and verified
programmatically).}
\label{tab:census}
\begin{tabularx}{\textwidth}{@{}lYY@{}}
\toprule
Object & Formula / origin & Value \\
\midrule
Group order & $|\mathrm{PSL}(2,7)| = 84(g-1)$ & 168 \\
Elements of orders 2 / 3 / 4 / 7 & census of $\mathrm{GL}(3,2)$ & 21 / 56 / 42 / 48 \\
Sylow numbers $n_2, n_3, n_7$ & subgroups $C_2, C_3, C_7$ & 21 / 28 / 8 \\
Vertices / edges / faces & $|G|/3,\ |G|/2,\ |G|/7$ & 56 / 84 / 24 \\
Flags (oriented / all) & $|G|,\ 2|G|$ & 168 / 336 \\
1-skeleton $\{7,3\}$ (Klein graph) & double coset $H_3\,t\,H_3$, $t^2 = e$ &
56 vertices, deg.~3, girth 7 \\
Face adjacency of $\{7,3\}$ & double coset $H_7\,t\,H_7$ & 24 vertices,
deg.~7, 84 edges \\
Euler characteristic & $V - E + F = 2 - 2g$ & $-4$ \\
Valence identity & $7F = 3V = 2E = |G|$ & 168 \\
\bottomrule
\end{tabularx}
\end{table}

\subsection{Provenance of the constants 224 and 672}

Both denominators of the DSI-4 observation are expressible through the census
of Table~\ref{tab:census} repeatedly and independently; the enumeration of
exact identities provides the criterion of ``census birth'': a constant is
not accidental if it lies in the intersection of many independent counting
families. For $224$ there are seven such appearances, for $672$ --- eight;
two cross-tower identities deserve note: $224 = 14(g{-}1)|_{g=17}$ is
the number of $C_3$-subgroups of the genus-17 Hurwitz triplet surface, and
$672 = 4V|_{g=7}$ is four times the vertex count of the Fricke--Macbeath
surface of genus 7.

\begin{table}[htbp]
\centering
\small
\caption{Provenance: all exact appearances of $224$ and $672$ in the census
of $\{7,3\}$ and the Hurwitz tower.}
\label{tab:prov}
\resizebox{\textwidth}{!}{%
\begin{tabular}{ll@{\hspace{2.2em}}ll}
\toprule
\multicolumn{2}{c}{$224$} & \multicolumn{2}{c}{$672$} \\
\cmidrule(r){1-2}\cmidrule(l){3-4}
$4\cdot 56$ & vertices of $\{7,3\}$: $4V$ & $4\cdot 168$ & $4|G|$ (4 spinor components) \\
$8\cdot 28$ & $8\,n_3$ ($C_3$-subgroups) & $2\cdot 336$ & $2|\mathrm{PGL}(2,7)|$ \\
$8\cdot 28$ & 28 bitangents $= 8N$ & $24\cdot 28$ & $F\cdot n_3$ (faces $\times$ $C_3$) \\
$32\cdot 7$ & $2^5\cdot 7$ (7 = order of $C$) & $12\cdot 56$ & $12V$ \\
$168+56$ & $|G| + V$ & $8\cdot 84$ & $8E$ \\
$2\cdot 84+56$ & $2E + V$ & $3\cdot 224$ & 3 = vertex valence \\
$14\cdot 16$ & $N(g{=}17)$: Hurwitz triplet & $4\cdot 168$ & $4V$ at genus 7 (Fricke--Macbeath) \\
\bottomrule
\end{tabular}%
}
\end{table}

\subsection{Uniqueness lemma: $9/224$ among census fractions}

Call a number \emph{census} if $q = 2^a 3^b 7^c$ with exponent budget
$a \le 6$, $b \le 1$, $c \le 1$ --- exactly the shape of every census
denominator (64 spin structures, 28 bitangents, 24 faces, 56 vertices,
84 edges, 168, 336, 672); denominators such as $25 = 5^2$, $49 = 7^2$,
$81 = 3^4$, $11^2$ do not occur in the census \emph{by construction}. An
exhaustive search over fractions $p/q$ with $q \le 5000$ from the semigroup
and $p \le 32$ gives: within $0.1\%$ of the measured $c_{K3}$ lies
\emph{exactly one} rational number --- $9/224$ (its unreduced writings
$18/448$, $27/672$ exhaust all hits; within $0.01\%$ the same unique
fraction). The continued fraction of the measured value,
$[0; 24, 1, 8, 19, \ldots]$, has exactly two census convergents: $1/24$
(deviation $+3.68\%$) and $9/224$ ($+0.0025\%$); the intermediate $1/25$ is
discarded by the filter as non-census. The exclusivity of $9/224$ is
therefore not a fitting artefact: any alternative of comparable accuracy
must have a denominator outside the $\{7,3\}$ census (Fig.~\ref{fig:scan}).

\begin{figure}[htbp]
\centering
\includegraphics[width=0.94\textwidth]{figs_en/fig224_scan.png}
\caption{The census filter: errors $|p/q - c_{K3}|$ for all census
denominators (blue) and a random sample of other fractions (grey). Inside
the $0.1\%$ band --- only $9/224 = 27/672$.}
\label{fig:scan}
\end{figure}

\subsection{The genus family $c(g)$ and the degeneracy theorem}

The main result of the attack: the observation $27/672$ admits \emph{two}
structural readings, indistinguishable on the Klein quartic but divergent on
the rest of the Hurwitz tower. Let a tower surface have genus $g$, group
order $|G| = 84(g-1)$, and the canonical $\{7,3\}$ tiling: then
$V = 28(g-1)$, $F = 12(g-1)$ and $N := n_3 = 14(g-1)$ is the number of
$C_3$-subgroups (vertex stabilisers). Branch~I reads $27/672$ as
$3^3/(4|G|)$ --- ``resonances per spinor slots over the group''; branch~II as
$(N-1)/(FN)$ --- ``all but one of the $C_3$-subgroups over
face$\times$subgroup slots'' (the $(N{-}1)/(24N)$ reading of Experiment 2).
Explicitly:
\begin{equation}\label{eq:branches}
c_I(g) = \frac{27}{4\cdot 84(g-1)} = \frac{9}{112(g-1)},
\qquad
c_{II}(g) = \frac{14(g-1) - 1}{12(g-1)\cdot 14(g-1)}.
\end{equation}

\begin{theorem}[degeneracy]\label{thm:degen}
The branches $c_I(g)$ and $c_{II}(g)$ coincide if and only if $g = 3$.
\end{theorem}
\begin{proof}
The equality $27\cdot 168(g-1)^2 = 336(g-1)\bigl(14(g-1)-1\bigr)$, after
cancelling $168(g-1)$, gives $27(g-1) = 28(g-1) - 2$, i.e.\ $g - 1 = 2$.
\end{proof}

A numerical check on the grid $g = 3..100$ confirms that the degeneracy
point is unique. Branch~III --- the framework hypothesis
$c_{III}(g) = b_{\mathrm{Ch}}(b_2(g))$ with hypothetical
$b_2 = 22, 46, 94, 190$ --- separates from both already at $g = 7$ and
behaves qualitatively differently: $b_{\mathrm{Ch}}(\lambda)
= 2\sin^2(\pi/\lambda) \sim 2\pi^2/\lambda^2$, i.e.\ $(g-1)^{-2}$ versus
$(g-1)^{-1}$. The prediction table:

\begin{table}[htbp]
\centering
\small
\caption{The genus family of the DSI amplitude $c(g)$: three branches on the
Hurwitz tower. At $g = 3$ all branches merge (quartic degeneracy); at
$g \ge 7$ they are distinguishable.}
\label{tab:genus}
\begin{tabular}{rrrrrrr}
\toprule
$g$ & $N = n_3$ & $c_I$ & $c_{II}$ & $c_{II}/c_I$ & $c_{III} = b_{\mathrm{Ch}}(b_2)$ &
$c_{III}/c_I$ \\
\midrule
3 & 28 & $9/224 = 0.040179$ & $9/224 = 0.040179$ & 1.0000 & 0.040509 & 1.0082 \\
7 & 84 & $3/224 = 0.013393$ & $83/6048 = 0.013724$ & 1.0247 & 0.009314 & 0.6954 \\
14 & 182 & $9/1456 = 0.006181$ & $181/28392 = 0.006375$ & 1.0313 & 0.002233 & 0.3613 \\
17 & 224 & $9/1792 = 0.005022$ & $223/43008 = 0.005185$ & 1.0324 & 0.000547 & 0.1089 \\
\bottomrule
\end{tabular}
\end{table}

\begin{findbox}[discriminator]
Branch~I decays as $(g-1)^{-1}$, branch~III as $(g-1)^{-2}$: at $g = 17$
they differ by a factor of nine. \emph{One} measurement of the DSI
modulation amplitude on any higher Hurwitz surface (nearest candidates:
Fricke--Macbeath $g = 7$ with $c_I = 3/224$, or the triplet $g = 17$)
suffices to select the mechanism. DSI-4 thus upgrades from a single
coincidence to a one-parameter falsifiable family (Fig.~\ref{fig:genus}).
\end{findbox}

\begin{figure}[htbp]
\centering
\includegraphics[width=0.94\textwidth]{figs_en/fig224_genus.png}
\caption{The genus family $c(g)$: branch~I $9/(4V)$, branch~II $(N-1)/(FN)$,
and branch~III $b_{\mathrm{Ch}}(b_2)$. The only degeneracy point
$c_I = c_{II}$ is $g = 3$ (Theorem~\ref{thm:degen}).}
\label{fig:genus}
\end{figure}

\subsection{The status of DSI-4 after the attack}

The upshot is threefold. The denominator $224 = 4\cdot 56$ is a census
constant of the $\{7,3\}$ tiling: it is expressible through the quartic
census in seven ways and has cross-tower twins ($N$ of the genus-17 triplet,
$4V$ at genus 7), making an accidental appearance highly unlikely. The
fraction $9/224$ is the unique census fraction near the measured value
(uniqueness lemma, exhaustive search). The numerator $9 = 3^2$ receives two
interpretations --- ``$3^3$ resonances over $4|G|$ spinor slots'' (branch~I)
and ``$N - 1$ out of $FN$ slots'' (branch~II) --- and only an experimental
tower can tell them apart: at genus 3 the mechanisms are degenerate by
Theorem~\ref{thm:degen}. What remains open is the dynamical derivation of
the numerator (why ``all but one'', or exactly $3^3$) --- the same
mechanism problem as for the $\gamma$ link; its statement is now precise.

% ==========================================================================
\section{The threshold $e \ge 1$: the ladder becomes a theorem}\label{sec:threshold}
% ==========================================================================

\subsection{The class of carrier-isometric pairs}

The Cuntz defect functional for a pair $X_1, X_2 \in M_n(\mathbb{C})$:
\begin{equation}\label{eq:F}
F(X_1, X_2) = \|X_1^*X_1 - I\|^2 + \|X_2^*X_2 - I\|^2
+ \|X_1^*X_2\|^2 + \|X_1X_1^* + X_2X_2^* - I\|^2,
\qquad e = F/n,
\end{equation}
with Frobenius norms; the exact Cuntz relations $X_i^*X_j = \delta_{ij}I$,
$X_1X_1^* + X_2X_2^* = I$ are realised in no $M_n$. The sharp trace theorem
of the audit gives the global floor $F \ge 5n/7$, with equality at the pair
$\sqrt{4/7}\,I$ --- a construction that \emph{contracts} the whole space.
Natural ``transports'' of the framework, however, do not contract the
carrier. We formalise this.

\begin{definition}[carrier-isometric pair]\label{def:carrier}
$X \in M_n(\mathbb{C})$ is a \emph{partial isometry} if all its nonzero
singular values equal 1 (equivalently: $XX^*X = X$; the carrier is the
initial subspace $(X^*X)\mathbb{C}^n$, on which $X$ is isometric). A pair
$(X_1, X_2)$ is \emph{carrier-isometric} if both operators are partial
isometries.
\end{definition}

The class of Definition~\ref{def:carrier} contains all structural models of
the framework: permutation matrices (unitary), truncated Cuntz-tree shifts
(partial isometries with projection $P_i = S_iS_i^*$), complementary
projections. Outside the class: $\sqrt{4/7}\,I$ (singular values
$\sqrt{4/7} \approx 0.756$) and weighted trees $c_iS_i$ with $c_i \ne 1$.

\subsection{Theorem A: the threshold}

\begin{theorem}[threshold $e \ge 1$]\label{thm:A}
For every carrier-isometric pair $(X_1, X_2) \in M_n(\mathbb{C})^2$:
\begin{equation}\label{eq:thmA}
F(X_1, X_2) \ge n, \qquad \text{i.e. } e \ge 1.
\end{equation}
Equality holds if and only if $P = X_1X_1^*$ and $Q = X_2X_2^*$ are
orthogonal ($PQ = 0$) and $p + q = n$, where $p = \mathrm{rank}\,P$,
$q = \mathrm{rank}\,Q$.
\end{theorem}

\begin{proof}
If $X_i$ are partial isometries, then $P = X_1X_1^*$, $Q = X_2X_2^*$ are
orthogonal projections. The audit identity $F = G(P,Q)$ with
$G(P,Q) = 2\,\mathrm{tr}P^2 + 2\,\mathrm{tr}Q^2 + 3\,\mathrm{tr}(PQ)
- 4\,\mathrm{tr}P - 4\,\mathrm{tr}Q + 3n$, together with
$\mathrm{tr}(P^2) = \mathrm{tr}(P)$ for projections, gives
\begin{equation}\label{eq:Fproj}
F = 3n - 2p - 2q + 3r, \qquad r = \mathrm{tr}(PQ).
\end{equation}
For projections the bounds $\max(0,\, p + q - n) \le r \le \min(p, q)$
hold: the lower one follows from $\mathrm{tr}(PQ) \ge \mathrm{tr}(P \wedge Q)
\ge p + q - n$ (dimension of the intersection of ranges), the upper one from
$\mathrm{tr}(PQ) = \mathrm{tr}(QPQ) \le \|P\|\,\mathrm{tr}(Q) = q$. If
$p + q \ge n$, then $F - n = 3r - 2(p+q-n) \ge 3(p+q-n) - 2(p+q-n) =
p+q-n \ge 0$; if $p + q \le n$, then $F - n = 2(n - p - q) + 3r \ge 0$. In
both cases equality forces $p + q = n$ and $r = 0$, i.e.\ $PQ = 0$.
\end{proof}

\begin{remark}[geometry of equality]\label{rem:equality}
The equality case is a pair of complementary orthogonal projections (e.g.\
$X_1 = \mathrm{diag}(I_p, 0)$, $X_2 = \mathrm{diag}(0, I_q)$, $p + q = n$):
the defect $F = n$ sits entirely in the ``fullness'' terms
$\|X_i^*X_i - I\|^2 = q$ and $p$, while the overlap and covering terms
vanish. Such a pair is degenerate as a Cuntz model: the operators are not
isometric on the whole space. The gauge group of $F$ is large ---
$(X_1, X_2) \mapsto (UX_1W,\, UX_2V)$ with arbitrary unitaries $U, V, W$
--- and all equality cases are equivalent to the block form above.
\end{remark}

\subsection{Theorem B: exact covering arithmetic}

\begin{theorem}[covering formula]\label{thm:B}
Let $X_1, X_2$ be partial isometries with $P + Q \le I$ (the combined ranges
do not cover the space) and $R := I - P - Q \ge 0$. Then
\begin{equation}\label{eq:thmB}
F = n + 2\,\mathrm{rank}(R) \quad \text{exactly.}
\end{equation}
\end{theorem}
\begin{proof}
From $P + Q \le I$ for projections it follows that $PQ = 0$. Counting the
four terms of~\eqref{eq:F} directly: $\|X_1^*X_1 - I\|^2 = \|P' - I\|^2$
where $P' = X_1^*X_1$ is the initial projection of rank $p$; since
$I - P = Q + R$, one has $\|P' - I\|^2 = \mathrm{tr}(I - P') = (n - p) =
q + \mathrm{rank}(R)$; symmetrically $\|X_2^*X_2 - I\|^2 = p +
\mathrm{rank}(R)$; next $\|X_1^*X_2\|^2 = \mathrm{tr}(PQ) = 0$ and
$\|P + Q - I\|^2 = \|R\|^2 = \mathrm{rank}(R)$. Summing and using
$p + q + \mathrm{rank}(R) = n$ yields $F = n + 2\,\mathrm{rank}(R)$.
\end{proof}

For the Cuntz-tree truncation at depth $k$ only the root is left uncovered:
$\mathrm{rank}(R) = 1$, and formula~\eqref{eq:thmB} instantly reproduces the
empirical finding of Experiment 1: $F = n + 2$, $e = 1 + 2/n \to 1$
(Fig.~\ref{fig:tree}; verified up to $k = 10$, $n = 2047$, accuracy
$10^{-13}$). The tree truncation is a \emph{near-extremal} of the class: it
pays exactly $2$ for the single uncovered dimension.

\begin{figure}[htbp]
\centering
\includegraphics[width=0.88\textwidth]{figs_en/fig7_tree.png}
\caption{Covering arithmetic: the Cuntz-tree truncation follows $F = n + 2$
exactly (Theorem~\ref{thm:B}, $\mathrm{rank}(R) = 1$); the class floor
$F = n$ is out of reach for the tree, and unitary pairs pay $2n$.}
\label{fig:tree}
\end{figure}

\subsection{Theorem C and the three-term defect arithmetic}

\begin{theorem}[unitary pairs]\label{thm:C}
If $X_1, X_2$ are unitary, then $F = 2n$ exactly, for every pair.
\end{theorem}
\begin{proof}
$P = Q = I$, the first two terms of~\eqref{eq:F} vanish,
$\|X_1^*X_2\|^2 = \mathrm{tr}(I) = n$ (a product of unitaries is unitary),
and the last term is $\|I\|^2 = n$.
\end{proof}

The permutation models of the framework are a special case: \emph{every}
unitary pair pays exactly $2n$, with the defect sitting entirely in the
\emph{overlap} of ranges ($r = n$). Formulas~\eqref{eq:Fproj}
and~\eqref{eq:thmB} provide the exact three-term arithmetic: the defect of a
carrier-isometric pair decomposes into (i) \emph{range overlap} with weight
$3r$, (ii) \emph{covering deficit} with weight $2\,\mathrm{rank}(R)$, and
(iii) the base unit $n$; the framework's ladder is exactly this arithmetic
evaluated on three families (Fig.~\ref{fig:ladder}): free scalars $5/7$
(outside the class), tree $1 + 2/n$, unitary pairs $2$.

\begin{figure}[htbp]
\centering
\includegraphics[width=0.94\textwidth]{figs_en/fig7_ladder.png}
\caption{The ladder as a theorem: on the class of carrier-isometric pairs
the floor is $e = 1$ (Theorem~\ref{thm:A}); the value $5/7$ is attainable
only by contracting the carrier (outside the class), and unitary pairs pay
exactly $e = 2$ (Theorem~\ref{thm:C}).}
\label{fig:ladder}
\end{figure}

\subsection{Consequence for the Sh1--Sh4 chain and the Sh.3(iii) bridge}

\begin{corollary}[class closure]\label{cor:D}
On the class of carrier-isometric pairs $\inf e = 1$; the value $5/7 < 1$
is unattainable. Hence no sequence of carrier-isometric models can drive
$e \to 0$: in the Sh1--Sh4 chain the step Sh2 strengthens to
$\eta \equiv 1$, and the Sh.3(iii) bridge (soficity $\Rightarrow$ matrix
models with $e \to 0$) is impossible for isometric-on-carrier transports ---
whatever permutation data one takes, the per-dimension defect stays at or
above 1.
\end{corollary}

\begin{remark}[the class boundary is optimal]
A weighted tree $c_iS_i$ leaves the class with a quadratic penalty
\begin{equation}\label{eq:weightedtree}
F = (n-1)[(c_1^2{-}1)^2 + (c_2^2{-}1)^2] + n + 2;
\end{equation}
at
$c^2 = 4/7$ this gives $e \approx 1.37$--$1.48$ (for $k = 3$--$7$) --- above
the class floor. Reaching $5/7$ requires contracting the \emph{entire}
carrier ($\sqrt{4/7}\,I$): the global minimum lives strictly outside the
class, so the threshold of Theorem~\ref{thm:A} is unavoidable within it.
\end{remark}

\subsection{Numerical verification}

All statements were verified by the experiment \texttt{exp7\_e1\_threshold.py}
(deterministic, fixed seeds); see Table~\ref{tab:verify}. Separately, the
supporting bounds $\max(0, p+q-n) \le \mathrm{tr}(PQ) \le \min(p,q)$ were
checked on $10^4$ random pairs of projections (0 violations), and the
identity~\eqref{eq:Fproj} to $5.9\cdot 10^{-16}$.

\begin{table}[htbp]
\centering
\small
\caption{Numerical verification of Theorems A/B/C (Experiment 7).}
\label{tab:verify}
\begin{tabularx}{\textwidth}{@{}lYY@{}}
\toprule
Check & Scope & Max deviation \\
\midrule
Identity $F = 3n - 2p - 2q + 3r$ & 24{,}000 random partial isometries,
$n \le 12$ & $5.9\cdot 10^{-16}$ \\
Threshold $e \ge 1$ (Theorem~A) & the same 24{,}000 pairs & 0 violations \\
Equality cases $F = n$ & 5 configurations $(n, p, q)$ & $8.9\cdot 10^{-16}$ \\
Covering formula (Theorem~B) & 500 random $(P, Q)$, $P + Q \le I$ &
$7.1\cdot 10^{-15}$ \\
Tree: $F = n + 2$, $\mathrm{rank}(R) = 1$ & $k = 2..10$ ($n$ up to 2047) &
$2.3\cdot 10^{-13}$ \\
Unitary pairs $F = 2n$ (Theorem~C) & 200 pairs, $n \le 32$ & $2.8\cdot 10^{-14}$ \\
Permutation pairs $F = 2n$ & 200 pairs & $1.4\cdot 10^{-14}$ \\
Bounds on $\mathrm{tr}(PQ)$ & 10{,}000 pairs of projections & 0 violations \\
\bottomrule
\end{tabularx}
\end{table}

% ==========================================================================
\section{Conclusions and next steps}\label{sec:conclusions}
% ==========================================================================

Assembled into one picture, the results of this session move the programme in
three ways. First, the ``obstruction ladder'' of the Cuntz defect is no
longer an empirical sequence of values but a theorem with exact arithmetic:
the threshold $e \ge 1$ on the class of carrier-isometric constructions is
proved (Theorem~\ref{thm:A}), the covering formula
$F = n + 2\,\mathrm{rank}(R)$ is exact (Theorem~\ref{thm:B}), unitary pairs
pay exactly $2n$ (Theorem~\ref{thm:C}), and Corollary~\ref{cor:D} closes the
Sh.3(iii) bridge for the whole class: sofic amplification cannot come from
isometric transports. Second, the DSI-4 observation acquired an analytical
environment: the constant $224 = 4\cdot 56$ is census-born (seven independent
identities, cross-tower twins), the fraction $9/224$ is unique among census
fractions (uniqueness lemma), and the genus family $c(g)$ with a single
degeneracy at $g = 3$ (Theorem~\ref{thm:degen}) makes the programme
falsifiable: the branches $(g-1)^{-1}$ and $(g-1)^{-2}$ differ by a factor
of $1.4$ already at $g = 7$ and by nine at $g = 17$. Third, the Choptuik
problem statement is fixed at both levels (physical and spectral), and its
solvability is assessed honestly: numerically --- solved; rigorously --- major
open problems (existence, uniqueness, stability of the DSS critical
solution); closed form --- the framework's programme with two precise gaps
(the mechanism and the $0.67\%$ offset).

Natural next steps: (a)~commit Theorems~\ref{thm:A}--\ref{thm:C} and
Corollary~\ref{cor:D} to the repository as version 3 of
\texttt{stability\_lemma} with tests on partial isometries;
(b)~set up a lattice/spectral experiment on the Fricke--Macbeath surface
($g = 7$): a measurement of the DSI amplitude would discriminate branches~I
and~III with a single number ($3/224$ versus $b_{\mathrm{Ch}}(46)$ --- a
factor of $1.44$); (c)~a high-precision computation of $\lambda_1$ of the
Klein quartic (Hejhal's algorithm for Maass forms) and of $\gamma$ by
Gundlach-type perturbation theory --- to test whether the monograph's
$\delta$-expansion can close the $0.67\%$ gap between $b_{\mathrm{Ch}}(7)$
and $\gamma$; (d)~analytic work on the numerator $9 = 3^2$: an attempt to
derive ``all but one'' ($N - 1$) from the cohomology of spin bundles over
$\{7,3\}$, or a proof that no such derivation exists --- which would itself
sharpen the honest attribution of DSI-4 as an observation.

\end{document}
